\subsection{First applications}

Let A be a Noetherian ring, V = Spec(A) its spectrum. In this number, we call a hypersurface in V any subset of the form \footnote{Recall that \(V(x)\) consists of the prime ideals of \(A\) containing \(x\).} V(x) with \(x \in A\) .

PROPOSITION 4. --- Let X be a closed subset of V, and let \(H_{1}, \ldots, H_{m}\) be hypersurfaces in V. Set \(X' = X \cap H_{1} \cap \ldots \cap H_{m}\) .

\begin{enumerate}
\def\labelenumi{\alph{enumi})}
\tightlist
\item
  For every closed subset Y of V contained in X′, we have
\end{enumerate}

\[
\operatorname{codim} (\mathrm{Y}, \mathrm{X} ^ {\prime}) \geqslant \operatorname{codim} (\mathrm{Y}, \mathrm{X}) - m.
\]

\begin{enumerate}
\def\labelenumi{\alph{enumi})}
\setcounter{enumi}{1}
\item
  We have codim(Z, X) ≤ m for every irreducible component Z of X'. If X' is nonempty, then codim(X', X) ≤ m.
\item
  If Z is a closed irreducible subset of V contained in X such that \(\text{codim}(Z, X) \leqslant m\), there exist hypersurfaces \(H_{1}^{\prime}, \ldots, H_{m}^{\prime}\) such that Z is an irreducible component of \(X \cap H_{1}^{\prime} \cap \ldots \cap H_{m}^{\prime}\) .
\end{enumerate}

Let \(\mathfrak{a}\) be an ideal of A and \(x_{1},\ldots,x_{m}\) elements of A such that \(\mathrm{X}=\mathrm{V}(\mathfrak{a})\) and \(\mathrm{H}_{i}=\mathrm{V}(x_{i})\) for \(1\leqslant i\leqslant m\) . Let Z be a closed irreducible subset of V contained in X; there exists a prime ideal p of A containing \(\mathfrak{a}\) and such that \(\mathrm{Z}=\mathrm{V}(\mathfrak{p})\) .

Suppose first that Z is contained in X' and denote by \(\xi_{i}\) the image of \(x_{i}\) in the Noetherian local ring \(\mathbf{B} = \mathbf{A}_{\mathfrak{p}} / \mathfrak{a}\mathbf{A}_{\mathfrak{p}}\) . By prop. 7, b) of § 1, no. 3, we have

\[
\operatorname{codim} (Z, X) = \dim (B), \quad \operatorname{codim} (Z, X ^ {\prime}) = \dim (B / (\xi_ {1} B + \dots + \xi_ {m} B)).
\]

By prop. 2 of no. 1, we therefore have

\[
\operatorname{codim} (Z, X ^ {\prime}) \geqslant \operatorname{codim} (Z, X) - m.\tag{11}
\]

If Z is an irreducible component of X', we have codim(Z, X') = 0, whence codim(Z, X) ≤ m; this proves b). We prove a) by taking, in the two members of (11), the lower bound over the set of irreducible components Z of Y.

Conversely, suppose we have \(\operatorname{codim}(\mathbf{Z},\mathbf{X})\leqslant m\), that is, \(\dim (\mathbf{B})\leqslant m\). Since every element of \(\mathbf{A}_{\mathfrak{p}}\) is the product of an invertible element of \(\mathbf{A}_{\mathfrak{p}}\) by the image of an element of A, the scholium of No. 2 proves the existence of elements \(x_1^{\prime},\ldots ,x_m^\prime\) of A whose images in B generate an ideal of definition of B. Set \(\mathrm{H}_i^{\prime} = \mathrm{V}(x_i^{\prime})\) for \(1\leqslant i\leqslant m\). It is clear that Z is an irreducible component of \(\mathbf{X}\cap \mathbf{H}_1^{\prime}\cap \dots \cap \mathbf{H}_m^{\prime}\) .

COROLLARY 1. --- Let H be a nonempty hypersurface in V. The codimension of H in V is equal to 0 or 1.

We have codim(H, V) = 1 if and only if H contains no irreducible component of V. If this is the case, then all irreducible components of H have codimension 1 in V.

For every irreducible component Z of H, we have

\[
0 \leqslant \operatorname{codim} (Z, V) \leqslant 1
\]

by prop. 4, b), and one has codim(Z, V) = 0 if and only if Z is an irreducible component of V. By definition,

\[
\operatorname{codim} (\mathrm{H}, \mathrm{V}) = \inf _ {\mathrm{Z}} \operatorname{codim} (\mathrm{Z}, \mathrm{V})
\]

where Z runs through the set of irreducible components of H. Cor. 1 follows immediately from these remarks.

COROLLARY 2. --- Let X be a closed irreducible subset of V and H a hypersurface in V. Only three cases are possible:

\begin{enumerate}
\def\labelenumi{\arabic{enumi})}
\item
  we have \(\mathbf{X}\subset \mathbf{H}\)
\item
  the set \(\mathbf{X} \cap \mathbf{H}\) is nonempty and each of its irreducible components \(Z\) satisfies \(\operatorname{codim}(\mathbf{Z}, \mathbf{X}) = 1\) ;
\item
  The set X ∩ H is empty.
\end{enumerate}

Suppose \(X' = X \cap H\) is nonempty and distinct from \(X\); every irreducible component \(Z\) of \(X'\) is distinct from \(X\) and satisfies codim \((Z, X) \leqslant 1\) by prop. 4, b). Cor. 2 follows immediately from this.

COROLLARY 3. --- If A is factorial (VII, § 3, no. 3, prop. 2), the prime ideals of height 1 of A are the principal ideals generated by the extremal elements of A. If moreover A is local, then dim(A/p) = dim(A) - 1 for every prime ideal p of height 1 of A.

Let \(x\) be an extremal element of A. Then \(Ax\) is a prime ideal because \(x\) is extremal, of height 1 because A is integral (no. 1, prop. 1). Let \(p\) be a prime ideal of height 1 of A. Then \(V(p)\) is an irreducible component of a hypersurface \(V(Ax)\) for a suitable \(x\) (prop. 4, c)). Let \(x = \prod_{i} y_{i}^{n_i}\) be a decomposition of \(x\) into products of extremal

elements such that \(y_{i}\) and \(y_{j}\) are coprime if \(i \neq j\). The irreducible components of \(V(Ax)\) are the \(V(Ay_{i})\). Therefore \(p = Ay_{i}\) for a suitable i. The last assertion follows from cor. 2 to prop. 2 of no. 1.

Remarks. --- 1) Suppose that A is a local, Noetherian, integral domain of dimension d; let x be a nonzero element of \(m_{A}\) and let \(H = V(x)\) . By cor. 2 of prop. 4, every irreducible component of H has codimension 1 in X, hence dimension \(\leqslant d - 1\) (§ 1, no. 2, prop. 3). By cor. 2 of prop. 2 of no. 1, H has dimension d - 1, and one of these components therefore has dimension d - 1; all of them do if A is catenary. However, in general there may exist an irreducible component of H of dimension \textless{} d - 1 (cf. p. 87, exercise 7).

\begin{enumerate}
\def\labelenumi{\arabic{enumi})}
\setcounter{enumi}{1}
\item
  Let \(x_1, \ldots, x_n\) be elements of A, and set \(\mathrm{H}_i = \mathrm{V}(x_i)\) for \(1 \leqslant i \leqslant n\). Suppose there exists a finitely generated A-module M with support \(\mathrm{V} = \operatorname{Spec}(\mathrm{A})\) such that \((x_1, \ldots, x_n)\) is a completely secant sequence for M. Then every irreducible component of \(\mathrm{H}_1 \cap \ldots \cap \mathrm{H}_n\) is of codimension \(n\) in V: this follows easily from the corollary of Prop. 3 of No. 2.
\item
  If the ideal a of the Noetherian ring A is generated by m elements, then ht(a) ≤ m. This follows immediately from prop. 4.
\end{enumerate}

PROPOSITION 5. --- Let A be a Noetherian ring and \(\mathfrak{p} \subset \mathfrak{q}\) a non-saturated chain of prime ideals of A. The set E of prime ideals \(\mathfrak{r}\) of A such that \(\mathfrak{p} \subset \mathfrak{r} \subset \mathfrak{q}\) is an infinite chain. We have \(\bigcup \mathfrak{r} = \mathfrak{q}\) and \(\bigcap \mathfrak{r} = \mathfrak{p}\) .

By replacing A with A/p, we reduce to the case where p = \{0\}.

By lemma 1 of Section 1, we have q = \(\bigcup_{r \in E} r\) and prop. 2 of II, § 1, no. 1 shows that E is infinite.

Let \(y \neq 0\) be an element of \(\bigcap_{r \in E} r\). The height of q is finite ( \(n^{\circ}1\) , cor. 1 of prop. 2), and one has \(ht(q) \geqslant 2\) by hypothesis. There therefore exists a prime ideal \(q' \subset q\) of height 2. The first part of the proof applied to \(q'\) shows that the set \(E'\) of prime ideals of height 1 contained in \(q'\) is infinite; each of these ideals contains y by hypothesis. Now the Noetherian local ring \(B = A_{q'} / yA_{q'}\) is of dimension 1 according to cor. 2 of prop. 2 of n° 1. For every \(r \in E'\) , the prime ideal \(r/yA_{q'}\) of B is therefore minimal; consequently, the Noetherian ring B has infinitely many minimal prime ideals, which is absurd (II, § 4, n° 3, cor. 3 of prop. 14). One therefore has \(\bigcap_{r \in E} r = \{0\}\) .

PROPOSITION 6.--- Let A be a Noetherian ring of dimension \(\geqslant 2\) , and \(h\) be an integer such that \(0 < h < \dim(A)\) .

\begin{enumerate}
\def\labelenumi{\alph{enumi})}
\item
  A possesses infinitely many prime ideals of height h.
\item
  If A is of finite dimension, it possesses infinitely many prime ideals p such that \(\operatorname{ht}(\mathfrak{p}) = h\) and \(\dim(\mathbf{A}/\mathfrak{p}) = \dim(\mathbf{A}) - h\) .
\end{enumerate}

Since the dimension of A is the upper bound of the heights of the prime ideals of A (§ 1, n° 3, prop. 8), since every prime ideal of A is of finite height (n° 1, cor. 1 of prop. 2) and since \(h < \dim(A)\) , there exists an integer \(n > h\) and a prime ideal \(p\) of height \(n\) , hence a chain \(p_0 \subset \ldots \subset p_n = p\) of prime ideals of length \(n\) . We have \(ht(p_i) = i\) for \(0 \leqslant i \leqslant n\) , whence \(ht(r) = h\) for every prime ideal \(r\) of A such that \(p_{h-1} \subset r \subset p_{h+1}\). The set E of these ideals is infinite according to prop. 5, whence \(a\) ).

If A is of finite dimension, one may suppose that one has \(n = \dim(A)\) in what precedes. For every prime ideal \(r \in E\) , one has \(\operatorname{ht}(r) = h\) , whence \(\dim(A/r) \leqslant n - h\) , and since \(r \subset p_{h+1} \subset \ldots \subset p_n\) is a chain of length \(n - h\) , one has \(\dim(A/r) = n - h\) .

There exist non-Noetherian integral domains of dimension 2 possessing only one prime ideal of height 1, for example a valuation ring of height 2 (VI, § 4, n° 4, prop. 5).

